% Sample-size / episode-count / cost derivation for resolving a per-memory
% causal effect via paired forced-in/forced-out episodes.
%
% Verified against this project's implementation:
%   - src/memory_ope/evaluation/power_formula.py (general reusable module)
%   - tests/test_power_formula.py (reproduces effect_size_analysis.json and
%     ground_truth_power_analysis.json exactly, and reproduces the 2x
%     discrepancy discussed in Remark 1 to confirm its source)
%   - src/memory_ope/evaluation/power_formula_validation.py (empirical
%     Monte Carlo check; results/power_formula_validation.json)
%   - src/memory_ope/evaluation/resolution_grid.py (the general grid over
%     delta, n_memories, rho; results/resolution_grid.json/.png)
%
% Paste directly into the paper; \(\S\)-numbers below are relative, not tied
% to any specific section number in the paper itself.

\subsection{Sample size for a paired forced-in/forced-out ground-truth design}
\label{sec:power-derivation}

\paragraph{Setup.} For a given memory $m$, the ground-truth procedure runs
$n$ \emph{paired} episodes: in pair $i$, the memory is forced \textsc{in}
for one episode (binary outcome $Y_{1,i}$) and forced \textsc{out} for the
other ($Y_{0,i}$), while the task instance, environment seed, and every
other memory's inclusion are held identical between the two arms of the
pair. This is exactly the design implemented by
\texttt{ground\_truth\_runner.py} (real ALFWorld) and
\texttt{\_common.paired\_design\_stats} (Stage~1 simulator).

Let $D_i = Y_{1,i} - Y_{0,i}$, so $\mathbb{E}[D_i] = \delta$, the memory's
true causal effect. We wish to choose $n$ (pairs per memory) large enough
to reject $H_0\!:\delta=0$ with probability (power) $1-\beta$ at
significance level $\alpha$, when the true effect is $\delta$.

\paragraph{Assumption 1 (common per-arm success probability).} Both arms'
marginal success probability is approximated by a single $p$:
$\Pr(Y_1=1)\approx\Pr(Y_0=1)\approx p$. This is exact when $\delta=0$ and a
good approximation whenever $|\delta|\ll p(1-p)$, which holds throughout
this project's real and planned settings ($p\in[0.5,0.85]$,
$\delta\in[0.01,0.20]$). Under this assumption,
\[
  \operatorname{Var}(Y_1) \approx \operatorname{Var}(Y_0) \approx p(1-p).
\]
(\emph{Consequence, checked empirically in \S\ref{sec:power-caveats}}: using
the exact arm-specific probabilities $p_1=p+\delta/2$, $p_0=p-\delta/2$
instead changes $\operatorname{Var}(D)$ by less than 1\% at every
$(p,\rho,\delta)$ combination this project uses — negligible relative to
the other approximations below.)

\paragraph{Assumption 2 (paired correlation $\rho$).} Define the
paired-outcome correlation
\[
  \rho \;=\; \frac{\operatorname{Cov}(Y_1,Y_0)}{\sqrt{\operatorname{Var}(Y_1)\operatorname{Var}(Y_0)}}.
\]
$\rho>0$ reflects shared context (task difficulty, which other memories were
co-included) driving both arms' success the same direction — exactly the
noise a paired design is built to cancel. $\rho$ is \emph{not} assumed
a priori in this project: it is estimated either from the Stage~1 simulator
(\texttt{\_common.paired\_design\_stats}, giving a planning proxy of
$\bar\rho\approx0.154$) or, once real paired episodes exist, directly and
empirically from them (\texttt{ground\_truth\_runner.paired\_correlation},
giving the much higher real ALFWorld value $\bar\rho\approx0.6355$ across the
four ground-truthed memories reported in this paper).

\paragraph{Variance of the paired difference.} By definition of covariance
and Assumption~1,
\[
  \operatorname{Var}(D)
  = \operatorname{Var}(Y_1) + \operatorname{Var}(Y_0) - 2\operatorname{Cov}(Y_1,Y_0)
  = 2p(1-p) - 2\rho\, p(1-p)
  \;=\; 2p(1-p)(1-\rho) \;=:\; \sigma_D^2.
\]

\paragraph{Assumption 3 (CLT / normal approximation).} For $n$ pairs, the
sample mean $\bar D_n = \frac{1}{n}\sum_i D_i$ satisfies, by the central
limit theorem,
\[
  \bar D_n \;\overset{\text{approx.}}{\sim}\; \mathcal N\!\left(\delta,\; \sigma_D^2/n\right).
\]
This is an approximation, not an exact statement: $D_i$ is a discrete random
variable taking values in $\{-1,0,1\}$, not continuous. \S\ref{sec:power-caveats}
reports how well this approximation holds empirically at the sample sizes
this project actually uses.

\paragraph{Sample-size formula.} Testing $H_0:\delta=0$ against
$H_1:\delta\neq0$ (two-sided) or $H_1:\delta>0$ / $\delta<0$ (one-sided —
used for this project's "detect a harmful memory" framing) at level $\alpha$
with power $1-\beta$, the standard \emph{one-sample} normal-approximation
sample-size formula is
\[
  \boxed{\,n_{\text{pairs}} \;=\; \frac{\left(z_{\alpha/2} + z_{\beta}\right)^2 \, \sigma_D^2}{\delta^2}\,}
  \qquad\text{(two-sided)}, \qquad
  n_{\text{pairs}} = \frac{\left(z_{\alpha} + z_{\beta}\right)^2 \, \sigma_D^2}{\delta^2}
  \quad\text{(one-sided)},
\]
where $z_{\alpha/2}=\Phi^{-1}(1-\alpha/2)$ and $z_\beta=\Phi^{-1}(1-\beta)$
(e.g.\ $z_{0.025}=1.96$, $z_{0.20}=0.8416$ for the conventional
$\alpha=0.05$, power$=0.80$ used throughout this project).

\paragraph{Episodes and cost.} Each pair costs two episodes (forced-in +
forced-out). For a store of $N$ independent memories, each ground-truthed to
the same $n_{\text{pairs}}$:
\[
  \text{total episodes} \;=\; 2 \, N \, n_{\text{pairs}},
  \qquad
  \text{cost} \;=\; (\text{total episodes}) \times (\text{cost per episode}).
\]

\subsection*{Remark 1: a common source of a 2$\times$ error}
\label{rem:factor-of-two}

A natural but \emph{incorrect} variant of this formula carries an extra
leading factor of 2:
\[
  n_{\text{pairs}}^{\text{(wrong)}} \;=\; \frac{2\left(z_{\alpha/2}+z_\beta\right)^2 \sigma_D^2}{\delta^2}
  \;=\; 2\, n_{\text{pairs}}.
\]
This is the correct formula for a \emph{different} design: an
\textsc{unpaired}, two-independent-sample comparison of two groups, each of
size $n$, where a single \emph{per-group} variance $\sigma^2$ (not yet
combined across groups) appears twice because
$\operatorname{Var}(\bar X_1 - \bar X_2) = \sigma^2/n + \sigma^2/n =
2\sigma^2/n$ for two independent samples. In our \emph{paired} design,
$\sigma_D^2$ is already $\operatorname{Var}(D)$ — it has already combined
both arms' variance and their correlation via the $2p(1-p)(1-\rho)$ term
derived above. Applying the two-sample formula's extra factor of 2 on top of
$\sigma_D^2$ double-counts that combination step and inflates
$n_{\text{pairs}}$ (hence total episodes and cost) by exactly $2\times$.
Concretely, at this project's real measured point ($p=0.6611$,
$\rho=0.6355$, $\delta=0.03$, $N=30$), the correct formula gives
$n_{\text{pairs}}\approx1424$ (85{,}463 total episodes, \$969); the
incorrect doubled formula gives $n_{\text{pairs}}\approx2849$ (170{,}927
episodes) — exactly the discrepancy caught, and resolved in favor of the
correct (undoubled) formula, while preparing this appendix.

\subsection{Empirical validation and caveats}
\label{sec:power-caveats}

The formula above was checked two ways, both zero-cost (pure simulation,
no API spend):

\begin{enumerate}
\item \textbf{Against this project's own prior implementations.} The
  general module reproduces \texttt{effect\_size\_analysis.json}'s
  resolution table (e.g.\ 1424.4 vs.\ the file's 1424.39 pairs/memory at
  $\delta=0.03$) and \texttt{ground\_truth\_power\_analysis.json}'s
  trade-off table to within floating-point rounding of their hardcoded
  $z$-values (1.96, 0.8416) vs.\ this module's full-precision $\Phi^{-1}$.
\item \textbf{Against Monte Carlo simulation of the actual test.} For seven
  $(p,\rho,\delta)$ settings spanning $p\in[0.30,0.80]$,
  $\rho\in[0,0.8]$, $|\delta|\in[0.03,0.30]$ (including this project's real
  ALFWorld point and a one-sided "detection" setting), $n_{\text{pairs}}$
  paired binary outcomes were drawn from an \emph{exact} bivariate Bernoulli
  distribution with the target marginals and correlation (not from this
  project's task-context-based simulators, which induce $\rho$ only
  indirectly), and a standard one-sample paired $t$-test was run on each of
  20{,}000 replications. \textbf{Empirical power landed in
  $[79.7\%,\,82.3\%]$ across all seven settings} — at or slightly above the
  nominal 80\% target, never meaningfully below it.
\end{enumerate}

The small ($\lesssim2.5$ percentage point), consistently
\emph{safe-direction} (more power than promised, not less) overshoot in
higher-$\rho$ settings was traced to $D_i$'s genuine discreteness ($D_i\in
\{-1,0,1\}$, not continuous): an exact noncentral-$t$ power calculation —
which still assumes normality of the individual $D_i$, only relaxing the
known-variance assumption of the $z$-test above — matches the nominal 80\%
almost exactly (79.3--79.8\% across the settings checked this way), meaning
the sample-size formula's own normal-approximation machinery is sound; the
residual gap comes specifically from how a $t$-test's sample-variance
estimator behaves on sparse $\{-1,0,1\}$-valued data (most pairs concordant,
$D_i=0$, at high $\rho$) rather than continuous Gaussian data.

\paragraph{Every assumption, listed.} For a reviewer or reader checking this
derivation: (i) a single shared $p$ approximates both arms' success
probability (Assumption~1); (ii) $\rho$ is treated as constant across the
effect-size range being planned for, estimated once from data available at
planning time; (iii) the CLT normal approximation for $\bar D_n$
(Assumption~3), whose finite-sample accuracy is empirically checked in
\S\ref{sec:power-caveats} rather than assumed; (iv) independence of pairs
within a memory and independence across memories (no shared randomness
between two different memories' ground-truth reruns); (v) a fixed,
conventional $\alpha=0.05$ and power $=0.80$, not derived from any
project-specific loss function; (vi) cost is linear in episode count at a
single blended cost-per-episode figure, itself an average over six
real-measured, unequally-costly task types (\S\ref{sec:setup}) — resolving
a store whose task-type mix differs from this project's would need a
re-weighted cost-per-episode, not this single constant.
